% ═══════════════════════════════════════════════════════════════════════
%  Stratégie 2 --- XAUUSD niveaux x10. Chiffres : macros et tables générées.
% ═══════════════════════════════════════════════════════════════════════
\section{Résultats in-sample}
\label{sec:x10_in_sample}

\lettrine[lines=2, findent=2pt, nindent=0pt]{C}{ette} section publie la
campagne complète sur 2019--2025. Elle est organisée pour répondre, dans
l'ordre, à quatre questions~: y a-t-il un réglage qui tient~? Y a-t-il une
période qui tient~? Y a-t-il une population de transactions qui tient~? Et le
contexte demandé par le mandat sert-il à quelque chose~?

La réponse est non aux quatre. Voici comment.

\subsection{La grille des vingt-sept configurations}

La grille explore trois axes, trois valeurs chacun~: la largeur de la zone de
surveillance, le seuil d'accélération et la distance du stop du
\emph{breakout}. C'est l'espace de configurations gelé avant la mesure, et
rien d'autre n'a été essayé.

% GÉNÉRÉ par scripts/build_x10_report_assets.py — ne pas éditer à la main.
% Toute valeur vient de results/xau_x10/*.json.
\begin{table}[htbp]
\centering
\footnotesize
\caption{Les vingt-sept configurations de la grille, classées par Sharpe net décroissant, au spread constant de sélection. L'étoile marque le centre géométrique, support des décompositions~: ce n'est pas une configuration retenue.}
\label{tab:x10_grid27}
\begin{tabular}{r r r r r r r r}
\toprule
$z$ & $a_{\min}$ & $k_s$ & \textbf{Sharpe net} & \textbf{Trades} & \textbf{E[R]} & \textbf{PF} & \textbf{Repli max} \\
\midrule
$0{,}5$ & $0{,}3$ & $1{,}50$ & $-1{,}178$ & $1\,020$ & $-0{,}1596$ & $0{,}760$ & $-44{,}0$\,\% \\
$0{,}5$ & $0{,}3$ & $1{,}00$ & $-1{,}226$ & $1\,070$ & $-0{,}1740$ & $0{,}746$ & $-47{,}7$\,\% \\
$0{,}5$ & $0{,}3$ & $0{,}75$ & $-1{,}251$ & $1\,093$ & $-0{,}1797$ & $0{,}743$ & $-49{,}3$\,\% \\
$0{,}5$ & $0{,}2$ & $1{,}50$ & $-1{,}331$ & $1\,299$ & $-0{,}1533$ & $0{,}759$ & $-52{,}0$\,\% \\
$0{,}5$ & $0{,}2$ & $0{,}75$ & $-1{,}344$ & $1\,372$ & $-0{,}1669$ & $0{,}751$ & $-56{,}6$\,\% \\
$0{,}5$ & $0{,}2$ & $1{,}00$ & $-1{,}381$ & $1\,351$ & $-0{,}1657$ & $0{,}745$ & $-55{,}2$\,\% \\
$1{,}0$ & $0{,}3$ & $1{,}50$ & $-1{,}394$ & $1\,333$ & $-0{,}1645$ & $0{,}757$ & $-51{,}3$\,\% \\
$1{,}0$ & $0{,}3$ & $0{,}75$ & $-1{,}431$ & $1\,432$ & $-0{,}1733$ & $0{,}751$ & $-56{,}8$\,\% \\
$1{,}0$ & $0{,}3$ & $1{,}00$ & $-1{,}469$ & $1\,397$ & $-0{,}1695$ & $0{,}744$ & $-55{,}3$\,\% \\
$1{,}5$ & $0{,}3$ & $0{,}75$ & $-1{,}538$ & $1\,618$ & $-0{,}1592$ & $0{,}755$ & $-60{,}3$\,\% \\
$1{,}5$ & $0{,}2$ & $0{,}75$ & $-1{,}555$ & $1\,930$ & $-0{,}1412$ & $0{,}772$ & $-65{,}2$\,\% \\
$1{,}0$ & $0{,}2$ & $1{,}50$ & $-1{,}598$ & $1\,653$ & $-0{,}1555$ & $0{,}752$ & $-61{,}4$\,\% \\
$1{,}0$ & $0{,}2$ & $0{,}75$ & $-1{,}607$ & $1\,760$ & $-0{,}1603$ & $0{,}748$ & $-65{,}0$\,\% \\
$1{,}0$ & $0{,}2$ & $1{,}00$ & $-1{,}661$ & $1\,723$ & $-0{,}1602$ & $0{,}737$ & $-63{,}7$\,\%~$\star$ \\
$1{,}5$ & $0{,}3$ & $1{,}00$ & $-1{,}666$ & $1\,579$ & $-0{,}1668$ & $0{,}737$ & $-60{,}7$\,\% \\
$1{,}5$ & $0{,}3$ & $1{,}50$ & $-1{,}668$ & $1\,498$ & $-0{,}1621$ & $0{,}736$ & $-58{,}0$\,\% \\
$1{,}5$ & $0{,}2$ & $1{,}50$ & $-1{,}669$ & $1\,809$ & $-0{,}1468$ & $0{,}757$ & $-64{,}2$\,\% \\
$1{,}5$ & $0{,}2$ & $1{,}00$ & $-1{,}707$ & $1\,891$ & $-0{,}1508$ & $0{,}749$ & $-65{,}2$\,\% \\
$1{,}0$ & $0{,}1$ & $0{,}75$ & $-1{,}735$ & $2\,058$ & $-0{,}1601$ & $0{,}740$ & $-70{,}3$\,\% \\
$0{,}5$ & $0{,}1$ & $1{,}50$ & $-1{,}849$ & $1\,607$ & $-0{,}1730$ & $0{,}699$ & $-65{,}1$\,\% \\
$1{,}0$ & $0{,}1$ & $1{,}00$ & $-1{,}855$ & $2\,017$ & $-0{,}1634$ & $0{,}723$ & $-70{,}2$\,\% \\
$1{,}0$ & $0{,}1$ & $1{,}50$ & $-1{,}861$ & $1\,948$ & $-0{,}1638$ & $0{,}733$ & $-68{,}0$\,\% \\
$0{,}5$ & $0{,}1$ & $0{,}75$ & $-1{,}884$ & $1\,678$ & $-0{,}1904$ & $0{,}694$ & $-69{,}9$\,\% \\
$0{,}5$ & $0{,}1$ & $1{,}00$ & $-1{,}905$ & $1\,655$ & $-0{,}1889$ & $0{,}689$ & $-68{,}3$\,\% \\
$1{,}5$ & $0{,}1$ & $0{,}75$ & $-2{,}009$ & $2\,222$ & $-0{,}1607$ & $0{,}735$ & $-75{,}0$\,\% \\
$1{,}5$ & $0{,}1$ & $1{,}50$ & $-2{,}196$ & $2\,092$ & $-0{,}1635$ & $0{,}720$ & $-74{,}0$\,\% \\
$1{,}5$ & $0{,}1$ & $1{,}00$ & $-2{,}220$ & $2\,172$ & $-0{,}1687$ & $0{,}709$ & $-75{,}8$\,\% \\
\bottomrule
\end{tabular}
\end{table}


\begin{figure}[htbp]
\centering
\includegraphics[width=\linewidth]{x10_fig03_heatmaps.png}
\caption{Carte du Sharpe net sur les trois axes. Chaque panneau fixe la
distance du stop et fait varier la largeur de zone (verticalement) et le seuil
d'accélération (horizontalement). L'échelle de couleur est commune aux trois
panneaux.}
\end{figure}

Trois lectures, qu'il faut faire sans les sur-interpréter.

\begin{description}
  \item[La grille est monotone selon le seuil d'accélération.] Les neuf
  configurations au seuil le plus bas occupent le bas du classement, les neuf
  au seuil le plus haut le haut. Un seuil plus exigeant arme moins souvent,
  donc trade moins, donc perd moins. \textbf{Ce n'est pas un signal de
  bord}~: c'est la signature d'un coût fixe par transaction qui domine tout le
  reste.
  \item[La largeur de zone agit dans le même sens, pour la même raison.] Une
  zone plus large arme plus souvent.
  \item[La distance du stop ne classe rien.] Elle n'agit que sur les
  \emph{breakouts} --- le stop du \emph{reversal} est une constante gelée ---
  ce qui rend ce tiers de la grille structurellement moins discriminant.
\end{description}

\subsection{Le plateau : aucune configuration retenue}

La règle de sélection gelée ne retient pas le meilleur point de la carte, mais
le \emph{centre d'un plateau}~: une configuration dont les voisins immédiats
tiennent aussi. Concrètement, il faut qu'au moins 60~\% des voisins aient un
Sharpe positif, et que la médiane des voisins atteigne la moitié du pic.

\begin{warningbox}
\textbf{\XTenPlateauPassing{} configuration sur vingt-sept passe ce test.} La
part de voisins à Sharpe positif vaut \XTenPlateauShareNeighbours{} partout,
puisque les vingt-sept Sharpes sont négatifs. La règle gelée prévoit
explicitement ce cas~: il n'y a pas de configuration retenue, et le verdict
est négatif. \textbf{On ne se rabat pas sur le pic}~--- un pic isolé dans une
grille de vingt-sept points est la signature d'un surajustement, pas d'un
réglage.
\end{warningbox}

Un détail arithmétique mérite d'être signalé, car il montre que la règle ne
s'est pas dégradée en passoire sur un échantillon perdant. Quand le pic est
négatif, « la moitié du pic » est un nombre \emph{plus élevé} que le pic
lui-même~: le seuil de médiane vaut \XTenPlateauThreshold{} pour un pic à
\XTenPeakSharpe. La seconde moitié du critère est donc devenue plus dure, et
elle échoue aussi --- la médiane des voisins du centre géométrique vaut
\XTenPlateauMedianNeighbour.

\paragraph{Plateau contre pic.} La table de décision demande de publier les
deux côte à côte. Le pic est \XTenPeakConfig, à \XTenPeakSharpe. La
configuration analysée --- le centre géométrique --- est \XTenAnalysedConfig,
à \XTenSharpeIS. L'écart vaut \XTenSharpeGapPeak{} de Sharpe. Il est réel, et
il n'a aucune conséquence~: les deux perdent. L'annexe~\ref{app:x10_plateau}
développe cette distinction.

\begin{insightbox}
\textbf{La configuration analysée dans tout ce qui suit n'est pas une
sélection.} Pour que les décompositions aient un support, il en faut une~;
le brief de campagne avait désigné \emph{avant la mesure} le centre
géométrique de la grille. Ce choix n'a pas été fait au vu des résultats, et il
ne doit jamais être présenté comme « la configuration de la stratégie ». Il
n'y en a pas.
\end{insightbox}

\subsection{Courbe d'équité et immersion}

\begin{figure}[htbp]
\centering
\includegraphics[width=\linewidth]{x10_fig04_equite.png}
\caption{Courbe d'équité par séance de la configuration analysée, base~1. La
perte est régulière~: pas de décrochage brutal, pas de période de
récupération.}
\end{figure}

\begin{figure}[htbp]
\centering
\includegraphics[width=\linewidth]{x10_fig05_immersion.png}
\caption{Courbe d'immersion --- l'écart au plus haut historique de la courbe
d'équité. Elle ne revient jamais durablement à zéro.}
\end{figure}

Le rendement total vaut \XTenTotalReturn{} sur les sept années, pour un repli
maximal de \XTenMaxDDPct. Le plus long épisode passé sous un plus haut
historique dure \XTenRecoveryDays{} séances sur \XTenSampleSessions~: la
courbe passe l'essentiel de sa vie sous son plus haut.

\subsection{Walk-forward annuel}

% GÉNÉRÉ par scripts/build_x10_report_assets.py — ne pas éditer à la main.
% Toute valeur vient de results/xau_x10/*.json.
\begin{table}[htbp]
\centering
\small
\caption{Walk-forward annuel de la configuration analysée. Aucune année n'est positive~; la perte s'atténue avec le temps sans jamais changer de signe.}
\label{tab:x10_walkforward}
\begin{tabular}{l r r r r r r}
\toprule
\textbf{Année} & \textbf{Trades} & \textbf{E[R]} & \textbf{PF} & \textbf{PnL} & \textbf{Réussite} & \textbf{Part du net} \\
\midrule
2019 & $149$ & $-0{,}0444$ & $0{,}965$ & $-114{,}2$~\$ & $33{,}6$\,\% & $1{,}8$\,\% \\
2020 & $246$ & $-0{,}2740$ & $0{,}580$ & $-2\,178{,}9$~\$ & $28{,}9$\,\% & $34{,}8$\,\% \\
2021 & $234$ & $-0{,}2196$ & $0{,}634$ & $-1\,465{,}2$~\$ & $27{,}4$\,\% & $23{,}4$\,\% \\
2022 & $247$ & $-0{,}2235$ & $0{,}660$ & $-1\,146{,}4$~\$ & $28{,}7$\,\% & $18{,}3$\,\% \\
2023 & $216$ & $-0{,}1479$ & $0{,}764$ & $-548{,}2$~\$ & $31{,}9$\,\% & $8{,}8$\,\% \\
2024 & $307$ & $-0{,}1096$ & $0{,}788$ & $-646{,}4$~\$ & $31{,}9$\,\% & $10{,}3$\,\% \\
2025 & $324$ & $-0{,}0920$ & $0{,}939$ & $-158{,}9$~\$ & $35{,}2$\,\% & $2{,}5$\,\% \\
\bottomrule
\end{tabular}
\end{table}


\begin{figure}[htbp]
\centering
\includegraphics[width=0.92\linewidth]{x10_fig06_esperance_annuelle.png}
\caption{Espérance nette par transaction, année par année. Le critère gelé
demandait qu'au moins 60~\% des années soient positives.}
\end{figure}

\textbf{\XTenPositiveYears{} année sur \XTenYearsCount{} est positive.} Le
critère en demandait 60~\%. En revanche, la seconde moitié du critère passe~:
aucune année ne porte plus de la moitié du résultat net, le maximum étant
\XTenMaxYearShare{}. La perte n'est donc pas l'accident d'un exercice.

La dernière année, à \XTenLastYearExpectancy{}~R, est moins négative que la
pire, \XTenWorstYear{} à \XTenWorstYearExpectancy{}~R, mais la trajectoire
n'est pas monotone. Deux effets mécaniques coexistent sans que cette campagne
les sépare~: un spread constant pèse moins en proportion quand le prix de l'or
monte, tandis que le garde-fou gain/risque rejette davantage de candidats en
fin d'échantillon. Aucun des deux n'établit une amélioration du signal.

\subsection{Décomposition par scénario et par séance}

% GÉNÉRÉ par scripts/build_x10_report_assets.py — ne pas éditer à la main.
% Toute valeur vient de results/xau_x10/*.json.
\begin{table}[htbp]
\centering
\small
\caption{Décomposition par séance, à l'heure de New York de la barre de décision. C'est la coupe la plus contrastée de la campagne.}
\label{tab:x10_sessions}
\begin{tabular}{l r r r r r}
\toprule
\textbf{Coupe} & \textbf{n} & \textbf{E[R]} & \textbf{PF} & \textbf{PnL} & \textbf{Réussite} \\
\midrule
Asie 18:00-03:00 & $595$ & $-0{,}2567$ & $0{,}602$ & $-3\,605{,}4$~\$ & $26{,}9$\,\% \\
Londres 03:00-08:00 & $379$ & $-0{,}1367$ & $0{,}841$ & $-910{,}6$~\$ & $28{,}2$\,\% \\
New York 08:00-17:00 & $749$ & $-0{,}0954$ & $0{,}807$ & $-1\,742{,}1$~\$ & $36{,}0$\,\% \\
\bottomrule
\end{tabular}
\end{table}


Les quatre scénarios perdent (section~\ref{sec:x10_scenarios}) et les trois
séances aussi. La campagne corrigée ne fait plus apparaître une séance qui
concentrerait à elle seule l'échec~: l'écart entre la meilleure et la pire
espérance n'est que de \metric{\XTenSessionSpreadR{}~R}. La plus faible est
\XTenWorstSessionName{} à \XTenWorstSessionR{}~R, tandis que la meilleure
reste négative à \XTenBestSessionR{}~R.

\begin{warningbox}
\textbf{Cette coupe ne suggère pas d'ajouter un filtre horaire.} Un
paramètre choisi après avoir vu la décomposition n'est pas un paramètre, c'est
une sélection --- et la table de décision l'interdit explicitement. Surtout,
aucune séance ne fournit ici de contraste assez large pour constituer une
hypothèse prioritaire de réouverture.
\end{warningbox}

\subsection{Décomposition par régime de volatilité}

% GÉNÉRÉ par scripts/build_x10_report_assets.py — ne pas éditer à la main.
% Toute valeur vient de results/xau_x10/*.json.
\begin{table}[htbp]
\centering
\small
\caption{Décomposition par régime de volatilité~: terciles d'ATR M5 en haut, tranches de pas de grille rapporté à l'ATR en bas. Aucune des quatre tranches n'est positive.}
\label{tab:x10_regimes}
\begin{tabular}{l r r r r r}
\toprule
\textbf{Coupe} & \textbf{n} & \textbf{E[R]} & \textbf{PF} & \textbf{PnL} & \textbf{Réussite} \\
\midrule
ATR bas & $575$ & $-0{,}1898$ & $0{,}725$ & $-2\,572{,}0$~\$ & $26{,}4$\,\% \\
ATR moyen & $574$ & $-0{,}1545$ & $0{,}730$ & $-2\,160{,}5$~\$ & $31{,}4$\,\% \\
ATR haut & $574$ & $-0{,}1362$ & $0{,}763$ & $-1\,525{,}6$~\$ & $35{,}7$\,\% \\
\midrule
$<3$ & $145$ & $-0{,}2645$ & $0{,}670$ & $-499{,}3$~\$ & $30{,}3$\,\% \\
$3$--$6$ & $700$ & $-0{,}1644$ & $0{,}693$ & $-2\,642{,}9$~\$ & $34{,}9$\,\% \\
$6$--$10$ & $560$ & $-0{,}0603$ & $0{,}866$ & $-1\,052{,}9$~\$ & $31{,}1$\,\% \\
$>10$ & $318$ & $-0{,}2793$ & $0{,}645$ & $-2\,063{,}0$~\$ & $23{,}6$\,\% \\
\bottomrule
\end{tabular}
\end{table}


La tranche intermédiaire --- celle où dix dollars valent entre six et dix
\code{ATR} --- est la moins mauvaise, à \XTenBestDollarBandR{}~R. Les deux
extrêmes sont les pires. C'est exactement la lecture annoncée par la note de
faisabilité~: quand le pas de grille vaut moins de trois \code{ATR}, le prix
traverse plusieurs niveaux par barre et la cible est du bruit~; quand il en
vaut plus de dix, la cible est hors de portée dans les 48~barres de durée de
vie. \textbf{Aucune des quatre tranches n'est positive.}

\subsection{Par raison de sortie : l'arithmétique complète de l'échec}

% GÉNÉRÉ par scripts/build_x10_report_assets.py — ne pas éditer à la main.
% Toute valeur vient de results/xau_x10/*.json.
\begin{table}[htbp]
\centering
\small
\caption{Répartition des sorties. Le stop rend exactement $-1$~R, ce qui vérifie le dimensionnement de bout en bout~; la cible rend un peu plus de deux R, et elle n'arrive pas assez souvent.}
\label{tab:x10_exits}
\begin{tabular}{l r r r r r}
\toprule
\textbf{Sortie} & \textbf{n} & \textbf{Part} & \textbf{E[R]} & \textbf{PF} & \textbf{PnL} \\
\midrule
STOP & $1\,072$ & $62{,}2$\,\% & $-1{,}0002$ & $0{,}000$ & $-23\,072{,}6$~\$ \\
TARGET (prochain x10) & $265$ & $15{,}4$\,\% & $2{,}0870$ & $\infty$ & $10\,980{,}1$~\$ \\
TIME (48 barres) & $279$ & $16{,}2$\,\% & $0{,}7515$ & $14{,}943$ & $5\,184{,}3$~\$ \\
SESSION (16:55 New York) & $107$ & $6{,}2$\,\% & $0{,}3127$ & $2{,}898$ & $650{,}0$~\$ \\
\bottomrule
\end{tabular}
\end{table}


\begin{figure}[htbp]
\centering
\includegraphics[width=\linewidth]{x10_fig09_raisons_sortie.png}
\caption{À gauche, combien de transactions sortent par chaque porte~; à
droite, ce que chaque porte rapporte en moyenne.}
\end{figure}

Cette table est le cœur du dossier, et elle est saine.

\begin{enumerate}
  \item Le stop rend \metric{\XTenStopR{}~R}, ce qui vérifie le
  dimensionnement de bout en bout.
  \item La cible rend \metric{\XTenTargetR{}~R} --- davantage que le rapport
  minimal de~1, parce qu'une transaction prise après un maintien ou un
  débordement a un stop plus serré que la géométrie nominale.
  \item Avec ces deux seuls chiffres, le point mort serait atteint pour un
  taux de cible de \metric{\XTenBreakevenTargetRate}. Le taux réalisé est de
  \metric{\XTenTargetShare}.
  \item Les sorties temporelles --- durée maximale et clôture de séance ---
  sont positives en moyenne, mais elles ne comblent pas l'écart.
\end{enumerate}

\begin{insightbox}
\textbf{Il faudrait plus que doubler le taux d'atteinte de la cible pour
seulement rentrer dans ses frais.} Ce n'est pas un écart qu'un réglage
rattrape. Un filtre qui produirait un tel gain sans réduire le nombre de
transactions ni dégrader le rapport gain/risque n'existe pas dans l'espace
exploré --- et les vingt-sept configurations de la grille le confirment.
\end{insightbox}

\subsection{Les coupes descriptives du mandat}

% GÉNÉRÉ par scripts/build_x10_report_assets.py — ne pas éditer à la main.
% Toute valeur vient de results/xau_x10/*.json.
\begin{table}[htbp]
\centering
\small
\caption{Les trois coupes descriptives du mandat~: extension VWAP, contexte dollar, retest. Aucune ne fait passer l'espérance du bon côté de zéro.}
\label{tab:x10_context_cuts}
\begin{tabular}{l r r r r r}
\toprule
\textbf{Coupe} & \textbf{n} & \textbf{E[R]} & \textbf{PF} & \textbf{PnL} & \textbf{Réussite} \\
\midrule
Reversal étendu au VWAP & $396$ & $-0{,}1991$ & $0{,}672$ & $-1\,673{,}0$~\$ & $30{,}8$\,\% \\
Reversal non étendu & $688$ & $-0{,}1509$ & $0{,}737$ & $-2\,538{,}4$~\$ & $29{,}7$\,\% \\
DXY favorable au trade & $933$ & $-0{,}1877$ & $0{,}711$ & $-5\,100{,}6$~\$ & $30{,}0$\,\% \\
DXY adverse au trade & $829$ & $-0{,}1455$ & $0{,}770$ & $-1\,621{,}2$~\$ & $32{,}1$\,\% \\
Avec retest du niveau & $769$ & $-0{,}1715$ & $0{,}701$ & $-3\,148{,}9$~\$ & $30{,}0$\,\% \\
Sans retest & $993$ & $-0{,}1650$ & $0{,}748$ & $-3\,572{,}9$~\$ & $31{,}7$\,\% \\
\bottomrule
\end{tabular}
\end{table}


\begin{description}
  \item[Extension au VWAP.] Les \emph{reversals} pris loin du \code{VWAP} de
  séance rendent \XTenExtendedR{}~R contre \XTenNotExtendedR{}~R pour les
  autres. Les deux branches sont négatives et la campagne corrigée ne soutient
  pas l'idée d'en faire un filtre bloquant. La spécification l'avait figée
  comme \emph{tracée et non bloquante}.
  \item[Contexte dollar.] Le contexte favorable rend
  \XTenDxyFavourableR{}~R contre \XTenDxyAdverseR{}~R pour le contexte adverse.
  Le garde-fou réduit la taille quand le dollar est adverse~; cette différence
  de poids doit être séparée de la comparaison en unités de risque.
  \item[Retest.] Le retest rend \XTenRetestR{}~R contre
  \XTenNoRetestR{}~R sans retest. L'écart est faible et les deux branches
  perdent~: la mesure ne soutient pas le retest comme confirmation.
\end{description}

\subsection{Ablations : retirer une composante à la fois}

Les sept combinaisons de retrait des trois filtres de contexte ont été
loguées comme essais \emph{avant} d'être mesurées.

% GÉNÉRÉ par scripts/build_x10_report_assets.py — ne pas éditer à la main.
% Toute valeur vient de results/xau_x10/*.json.
\begin{table}[htbp]
\centering
\small
\caption{Ablations des trois filtres de contexte, à la configuration analysée. Le classement est monotone~: plus on retire de contexte, plus on perd.}
\label{tab:x10_ablations}
\begin{tabular}{c c l r r r r r}
\toprule
\textbf{EMA50 H1} & \textbf{VWAP} & \textbf{DXY} & \textbf{Trades} & \textbf{Sharpe} & \textbf{E[R]} & \textbf{PF} & \textbf{PnL} \\
\midrule
\checkmark & \checkmark & \checkmark \emph{(spec)} & $1\,762$ & $-1{,}941$ & $-0{,}1679$ & $0{,}728$ & $-6\,721{,}8$~\$ \\
\checkmark & \checkmark & --- & $1\,764$ & $-2{,}001$ & $-0{,}1688$ & $0{,}741$ & $-7\,579{,}0$~\$ \\
\checkmark & --- & \checkmark & $1\,832$ & $-1{,}948$ & $-0{,}1716$ & $0{,}726$ & $-6\,837{,}1$~\$ \\
\checkmark & --- & --- & $1\,838$ & $-2{,}088$ & $-0{,}1698$ & $0{,}732$ & $-7\,757{,}9$~\$ \\
--- & \checkmark & \checkmark & $1\,942$ & $-2{,}118$ & $-0{,}1722$ & $0{,}727$ & $-7\,148{,}9$~\$ \\
--- & \checkmark & --- & $1\,956$ & $-2{,}168$ & $-0{,}1705$ & $0{,}737$ & $-7\,973{,}8$~\$ \\
--- & --- & \checkmark & $2\,146$ & $-2{,}149$ & $-0{,}1717$ & $0{,}730$ & $-7\,372{,}7$~\$ \\
--- & --- & --- & $2\,168$ & $-2{,}209$ & $-0{,}1699$ & $0{,}742$ & $-8\,199{,}2$~\$ \\
\bottomrule
\end{tabular}
\end{table}


\textbf{Les trois filtres du mandat font ce qu'on leur demande, et cela ne
suffit pas.} Chaque retrait dégrade le Sharpe, et le retrait des trois le fait
passer de \XTenSharpeIS{} à \XTenAblationAllOffSharpe. Le classement est
parfaitement monotone.

Deux précisions comptent.

\begin{itemize}
  \item Retirer le filtre dollar ne change la population que de
  \XTenDxyTradeDelta{} transactions, mais dégrade le Sharpe et accroît la perte
  de \XTenDxySavings. C'est cohérent avec son rôle principal~: réduire la
  taille plutôt que bloquer les signaux.
  \item Retirer la moyenne horaire ouvre plus de transactions que retirer le
  \code{VWAP}. Les deux moitiés du filtre de contexte des \emph{breakouts} ne
  mordent donc pas également.
\end{itemize}

\begin{insightbox}
\textbf{On ne redresse pas une espérance franchement négative en ajoutant des
filtres qui la rapprochent de zéro au millième près.} C'est exactement ce que
montre la colonne des espérances de la table ci-dessus~: elle bouge à peine
d'une ligne à l'autre, alors que le Sharpe et le montant perdu, eux, bougent
beaucoup. Le contexte est utile --- il corrige un signal dont l'espérance
brute est trop négative pour être rattrapée.
\end{insightbox}

\subsection{Concentration du résultat}

% GÉNÉRÉ par scripts/build_x10_report_assets.py — ne pas éditer à la main.
% Toute valeur vient de results/xau_x10/*.json.
\begin{table}[htbp]
\centering
\small
\caption{Concentration du résultat. Les parts sont rapportées au profit \emph{brut}~: une part d'un net négatif n'aurait pas de sens.}
\label{tab:x10_concentration}
\begin{tabular}{l r r}
\toprule
\textbf{Concentration} & \textbf{Montant} & \textbf{Part du profit brut} \\
\midrule
Profit brut cumulé & $17\,973{,}64$~\$ & --- \\
Cinq meilleures séances & $1\,078{,}67$~\$ & $6{,}0$\,\% \\
Dix meilleurs trades & $1\,776{,}65$~\$ & $9{,}9$\,\% \\
Meilleur mois (2019-10) & $401{,}95$~\$ & --- \\
\midrule
Séances gagnantes & $397$ & --- \\
Séances perdantes & $766$ & --- \\
\bottomrule
\end{tabular}
\end{table}


Les parts du résultat \emph{net} ne sont pas publiables ici~: le net est
négatif, et « tant pour cent d'un nombre négatif » n'a pas de sens. Les
montants et les parts du profit \emph{brut} le sont.

Le constat est qu'\textbf{il n'y a aucune concentration pathologique}. Les
cinq meilleures séances représentent \XTenTopFiveDaysShare{} du profit brut,
les dix meilleures transactions \XTenTopTenTradesShare. Le meilleur mois de
sept années rapporte \XTenBestMonthPnl. Et le compte des séances donne
\XTenProfitableSessions{} gagnantes contre \XTenLosingSessions{} perdantes.

\begin{warningbox}
\textbf{L'absence de concentration rend la perte d'autant plus crédible.} Sur
une stratégie gagnante, une performance portée par cinq journées serait un
signal d'alarme. Ici c'est l'inverse~: la perte est diffuse, étalée sur sept
ans et sur des centaines de séances. Il n'y a pas d'événement singulier à
invoquer, pas de période exceptionnelle à écarter, rien qui permette de dire
« sans ceci, cela aurait marché ».
\end{warningbox}
